\documentclass[10pt,aspectratio=169]{beamer}
% Preámbulo modular para presentaciones Beamer (Modelación Estadística)
% Diseño y paleta institucional del Tecnológico de Monterrey

\documentclass[aspectratio=169,xcolor={dvipsnames,table}]{beamer}

\usepackage[utf8]{inputenc}
\usepackage[spanish,mexico]{babel}
\usepackage{amsmath,amssymb,amsfonts,mathtools}
\usepackage{graphicx}
\usepackage{listings}
\usepackage{booktabs}
\usepackage{tikz}
\usetikzlibrary{matrix,arrows,decorations.pathmorphing,shapes.geometric,positioning}

% Paleta Institucional Tecnológico de Monterrey
\definecolor{TecAzul}{HTML}{0039A6}       % Azul TEC: color primario institucional
\definecolor{TecAzulOscuro}{HTML}{1A2E51} % Azul marino oscuro
\definecolor{TecRojo}{HTML}{EC2661}       % Rojo de acento (paleta miTec)
\definecolor{TecGris}{HTML}{646464}       % Gris neutro para comentarios y subtítulos
\definecolor{TecFondoBloque}{HTML}{F4F6F9}% Fondo suave para bloques

% Configuración de Tema Beamer
\usetheme{Boadilla}
\usecolortheme[named=TecAzulOscuro]{structure}
\setbeamercolor{alerted text}{fg=TecRojo}
\setbeamercolor{palette primary}{bg=TecAzulOscuro,fg=white}
\setbeamercolor{palette secondary}{bg=TecAzul,fg=white}
\setbeamercolor{palette tertiary}{bg=TecAzulOscuro!80!black,fg=white}
\setbeamercolor{title}{fg=TecAzulOscuro,bg=TecFondoBloque}
\setbeamercolor{frametitle}{fg=TecAzulOscuro,bg=TecFondoBloque}

% Bloques personalizados elegantes
\setbeamercolor{block title}{bg=TecAzulOscuro,fg=white}
\setbeamercolor{block body}{bg=TecFondoBloque,fg=black}
\setbeamercolor{block title alerted}{bg=TecRojo,fg=white}
\setbeamercolor{block body alerted}{bg=TecRojo!8,fg=black}
\setbeamercolor{block title example}{bg=TecAzul,fg=white}
\setbeamercolor{block body example}{bg=TecAzul!8,fg=black}

% Redondear bordes y añadir sombra sutil
\setbeamertemplate{blocks}[rounded][shadow=true]
\setbeamertemplate{navigation symbols}{}

% Pie de página limpio con número de diapositiva
\setbeamertemplate{footline}{
  \leavevmode%
  \hbox{%
  \begin{beamercolorbox}[wd=.7\paperwidth,ht=2.25ex,dp=1ex,left]{author in head/foot}%
    \hspace*{2ex}\insertshortauthor~~(\insertshortinstitute) \hfill \insertshorttitle
  \end{beamercolorbox}%
  \begin{beamercolorbox}[wd=.3\paperwidth,ht=2.25ex,dp=1ex,right]{date in head/foot}%
    \insertshortdate{}\hspace*{2em}
    \insertframenumber{} / \inserttotalframenumber\hspace*{2ex}
  \end{beamercolorbox}}%
  \vskip0pt%
}

% Configuración de Listings de Python para visualización pedagógica de código
\lstset{
    language=Python,
    basicstyle=\ttfamily\footnotesize,
    keywordstyle=\color{TecAzulOscuro}\bfseries,
    stringstyle=\color{TecRojo},
    commentstyle=\color{TecGris}\itshape,
    showstringspaces=false,
    breaklines=true,
    frame=lines,
    rulecolor=\color{TecAzulOscuro!30},
    backgroundcolor=\color{TecFondoBloque},
    aboveskip=0.5em,
    belowskip=0.5em,
    literate={á}{{\'a}}1 {é}{{\'e}}1 {í}{{\'i}}1 {ó}{{\'o}}1 {ú}{{\'u}}1
             {Á}{{\'A}}1 {É}{{\'E}}1 {Í}{{\'I}}1 {Ó}{{\'O}}1 {Ú}{{\'U}}1
             {ñ}{{\~n}}1 {Ñ}{{\~N}}1 {¿}{{?`}}1 {¡}{{!`}}1
}

% Metadatos institucionales por defecto
\title[Modelación Estadística]{Modelación Estadística para Ciencia de Datos e Ingeniería}
\author[J. Castillo Colmenares]{Juliho David Castillo Colmenares}
\institute[Tec de Monterrey]{Tecnológico de Monterrey \\ Escuela de Ingeniería y Ciencias}
\date{\today}

% Comandos y macros estadísticas compartidas para las presentaciones Beamer (Metropolis)

% Conjuntos numéricos básicos
\newcommand{\R}{\mathbb{R}}
\newcommand{\N}{\mathbb{N}}
\newcommand{\Q}{\mathbb{Q}}
\newcommand{\Z}{\mathbb{Z}}
\newcommand{\set}[1]{\left\{ #1 \right\}}
\newcommand{\abs}[1]{\left|#1\right|}
\newcommand{\norm}[1]{\left\| #1 \right\|}

% Comandos estadísticos y probabilísticos del curso
\newcommand{\Var}{\operatorname{Var}}
\newcommand{\cov}{\operatorname{Cov}}
\newcommand{\corr}{\rho}
\newcommand{\comb}[2]{\begin{pmatrix} #1 \\ #2 \end{pmatrix}}
\newcommand{\s}{\sigma}
\newcommand{\particion}{\mathfrak{P}}
\newcommand{\card}[1]{\#\left( #1 \right)}
\newcommand{\seq}[3]{\set{#1}_{#2}^{#3}}
\newcommand{\E}{\mathbb{E}}
\newcommand{\Prob}{\mathbb{P}}

% Entornos pedagógicos y cajas en español académico natural
\newenvironment{discusion}{%
  \begin{alertblock}{Pregunta de Discusión}%
}{%
  \end{alertblock}%
}

\newenvironment{reflexion}{%
  \begin{alertblock}{Reflexión para el Aula}%
}{%
  \end{alertblock}%
}

\newenvironment{paradoja}{%
  \begin{alertblock}{Paradoja de Exploración}%
}{%
  \end{alertblock}%
}

\newenvironment{motivacion}{%
  \begin{exampleblock}{Motivación: Ciencia de Datos en la Práctica}%
}{%
  \end{exampleblock}%
}

\newenvironment{retoclase}{%
  \begin{block}{Reto Rápido en Equipo}%
}{%
  \end{block}%
}


\title[Estimación de proporciones]{Sección 6.5: Estimación de una proporción y diferencia entre dos proporciones}
\subtitle{Wald, Wilson y comparación de tasas}
\author[J. Castillo Colmenares]{Juliho Castillo Colmenares}
\institute{Tecnológico de Monterrey}
\date{}

\begin{document}
\begin{frame}[plain]\vspace{-0.8cm}\titlepage\end{frame}

\begin{frame}{Hoja de ruta --- capítulo 6}
  \footnotesize
  \begin{itemize}\itemsep=0.08cm
    \item \textbf{06.02} Intervalos para medias
    \item \textbf{06.04} Errores estándar
    \item \textbf{06.05} Proporciones y diferencias \emph{(hoy)}
    \item \textbf{06.07} Tamaño muestral
  \end{itemize}
  \begin{block}{Objetivo didáctico}
    Construir intervalos para una proporción y para la diferencia de dos
    proporciones, distinguiendo Wald de Wilson.
  \end{block}
\end{frame}

\begin{frame}{Motivación: respuestas categóricas}
  \footnotesize
  Las proporciones describen éxitos Bernoulli: tasas de clics, preferencias y
  métricas de clasificación.
  \begin{alertblock}{Pregunta guía}
    ¿Cómo cuantificar la incertidumbre de $\hat p=X/n$ y de una diferencia de
    tasas?
  \end{alertblock}
\end{frame}

\begin{frame}{Intervalo de Wald}
  \footnotesize
  Si las condiciones de muestra grande se cumplen,
  \[
    \hat p\pm Z_{\alpha/2}\sqrt{\frac{\hat p(1-\hat p)}{n}}.
  \]
  Se basa en la aproximación normal de la distribución muestral de $\hat p$.
\end{frame}

\begin{frame}{Limitación de Wald}
  \footnotesize
  Para muestras pequeñas o proporciones cercanas a $0$ o $1$, Wald puede tener
  cobertura real inferior a la nominal y producir límites fuera de $[0,1]$.
  \begin{block}{Condiciones de referencia}
    La nota usa $n\hat p\ge5$ y $n(1-\hat p)\ge5$ (o, de forma conservadora,
    $\ge10$).
  \end{block}
\end{frame}

\begin{frame}{Intervalo de Wilson}
  \footnotesize
  Invirtiendo directamente la prueba de score se obtiene
  \[
  \frac{1}{1+Z^2/n}\left[\hat p+\frac{Z^2}{2n}\pm
  Z\sqrt{\frac{\hat p(1-\hat p)}{n}+\frac{Z^2}{4n^2}}\right].
  \]
  El intervalo no se centra necesariamente en $\hat p$.
\end{frame}

\begin{frame}{Diferencia de dos proporciones}
  \footnotesize
  Para muestras independientes suficientemente grandes,
  \[
    (\hat p_1-\hat p_2)\pm Z_{\alpha/2}
    \sqrt{\frac{\hat p_1(1-\hat p_1)}{n_1}+
    \frac{\hat p_2(1-\hat p_2)}{n_2}}.
  \]
  \begin{alertblock}{Lectura}
    El estimador central es la diferencia y la varianza combina ambos grupos.
  \end{alertblock}
\end{frame}

\begin{frame}{Condiciones y elección del método}
  \footnotesize
  Wald es una aproximación de muestra grande; Wilson es preferible cuando los
  conteos son pequeños o la proporción está cerca de un extremo.
  \begin{alertblock}{Regla}
    Primero se revisan las condiciones; después se elige la fórmula.
  \end{alertblock}
\end{frame}

\begin{frame}{Puente computacional 1/4: estimador}
  \footnotesize
  \[
    \hat p=\frac{X}{n},\qquad
    \operatorname{Var}(\hat p)\approx\frac{p(1-p)}{n}.
  \]
  La estimación reemplaza $p$ por $\hat p$ en el error estándar de Wald.
\end{frame}

\begin{frame}{Puente computacional 2/4: corrección Wilson}
  \footnotesize
  Wilson conserva $p$ en el denominador durante la inversión del score y
  resuelve la desigualdad resultante.
  \begin{block}{Consecuencia}
    El centro y la amplitud se ajustan automáticamente cerca de los límites.
  \end{block}
\end{frame}

\begin{frame}{Puente computacional 3/4: dos grupos}
  \footnotesize
  \[
    \widehat{SE}(\hat p_1-\hat p_2)=
    \sqrt{\frac{\hat p_1(1-\hat p_1)}{n_1}+
    \frac{\hat p_2(1-\hat p_2)}{n_2}}.
  \]
  Cada muestra aporta su propia incertidumbre.
\end{frame}

\begin{frame}{Puente computacional 4/4: límites}
  \footnotesize
  \begin{itemize}\itemsep=0.12cm
    \item Wald puede exceder $[0,1]$.
    \item Wilson mantiene límites plausibles.
    \item La diferencia $p_1-p_2$ vive en $[-1,1]$.
  \end{itemize}
\end{frame}

\begin{frame}{Aplicación de lectura 1/4 --- Wald}
  \footnotesize
  La nota formula el IC normal para $p$ cuando se cumplen las condiciones de
  muestra grande.
  \begin{block}{Planteamiento}
    Identificar $X$, $n$, $\hat p$ y el valor crítico.
  \end{block}
\end{frame}

\begin{frame}{Aplicación de lectura 1/4 --- solución}
  \footnotesize
  \[
    IC_W=\hat p\pm Z_{\alpha/2}
    \sqrt{\frac{\hat p(1-\hat p)}{n}}.
  \]
  La aproximación es razonable cuando éxitos y fracasos esperados son
  suficientemente numerosos.
\end{frame}

\begin{frame}{Aplicación de lectura 2/4 --- Wilson}
  \footnotesize
  Para una muestra pequeña o una proporción extrema, la nota recomienda
  invertir el score en lugar de usar Wald.
  \begin{block}{Planteamiento}
    Sustituir la fórmula simétrica por la expresión de Wilson.
  \end{block}
\end{frame}

\begin{frame}{Aplicación de lectura 2/4 --- solución}
  \footnotesize
  El factor $1/(1+Z^2/n)$ y el término $Z^2/(4n^2)$ corrigen el centro y la
  amplitud.
  \begin{alertblock}{Resultado}
    Wilson mantiene una cobertura real más cercana al nivel nominal.
  \end{alertblock}
\end{frame}

\begin{frame}{Aplicación de lectura 3/4 --- dos proporciones}
  \footnotesize
  Dos poblaciones independientes tienen tasas $p_1$ y $p_2$ y tamaños $n_1$ y
  $n_2$.
  \begin{block}{Planteamiento}
    Estimar $p_1-p_2$ mediante $\hat p_1-\hat p_2$.
  \end{block}
\end{frame}

\begin{frame}{Aplicación de lectura 3/4 --- solución}
  \footnotesize
  \[
    IC=(\hat p_1-\hat p_2)\pm Z_{\alpha/2}\widehat{SE}.
  \]
  El intervalo incluye cero cuando la diferencia observada es compatible con
  ausencia de diferencia.
\end{frame}

\begin{frame}{Aplicación de lectura 4/4 --- condiciones}
  \footnotesize
  La validez de la aproximación normal depende de los conteos esperados
  $n_i\hat p_i$ y $n_i(1-\hat p_i)$.
  \begin{block}{Planteamiento}
    Revisar las condiciones antes de reportar un intervalo.
  \end{block}
\end{frame}

\begin{frame}{Aplicación de lectura 4/4 --- decisión}
  \footnotesize
  Si las condiciones fallan, se prefiere Wilson (o Agresti--Coull) para una
  proporción; para dos tasas se requiere una aproximación adecuada al diseño.
  \begin{exampleblock}{Resultado}
    El método se elige por la información disponible, no solo por conveniencia.
  \end{exampleblock}
\end{frame}

\begin{frame}{Síntesis de la sección}
  \footnotesize
  \begin{itemize}\itemsep=0.12cm
    \item Wald es simple y exige condiciones de muestra grande.
    \item Wilson corrige centro, amplitud y cobertura.
    \item Dos proporciones combinan dos errores estándar.
    \item Los conteos esperados determinan la aproximación apropiada.
  \end{itemize}
\end{frame}

\begin{frame}{Transición curricular}
  \footnotesize
  \begin{block}{Lo que queda establecido}
    Las proporciones requieren respetar el soporte $[0,1]$ y las condiciones de
    la aproximación elegida.
  \end{block}
  \begin{alertblock}{Siguiente sección}
    \textbf{06.07 Tamaño de una muestra}: elegir $n$ antes de observar los
    datos, controlando confianza y margen de error.
  \end{alertblock}
\end{frame}
\end{document}
